% secciones/01_introduccion.tex
\section{Introducción}
La integración de robots móviles en entornos hospitalarios ha pasado de ser una visión futurista a una necesidad operativa crítica. En el escenario actual, los sistemas sanitarios se enfrentan a desafíos sin precedentes: una población envejecida que demanda más servicios, una escasez global de personal de enfermería y la imperante necesidad de extremar los protocolos de asepsia para evitar infecciones nosocomiales. Estudios recientes en gestión hospitalaria estiman que el personal de enfermería dedica entre un 20\% y un 30\% de su jornada laboral a tareas logísticas de bajo valor clínico, como el transporte manual de muestras al laboratorio o la búsqueda de suministros en almacenes centrales. Esta ineficiencia no solo incrementa el agotamiento del personal (burnout), sino que reduce drásticamente el tiempo disponible para el cuidado directo y la supervisión del paciente, factores que impactan directamente en la calidad asistencial y en la tasa de errores médicos.

En este contexto, la automatización de la logística interna surge como una solución tecnológica transformadora. Un sistema robótico autónomo no es simplemente un vehículo de transporte; es un nodo inteligente en una red asistencial que permite una trazabilidad total de los suministros y una optimización de los flujos de trabajo. Al delegar el transporte de muestras biológicas, medicamentos controlados y suministros estériles en una plataforma autónoma, el hospital puede garantizar un flujo constante de materiales ("just-in-time"), eliminando cuellos de botella y minimizando el riesgo de contaminación cruzada derivado del tránsito humano excesivo entre zonas críticas.

El objetivo principal de este trabajo es el diseño integral de un robot mensajero hospitalario que no solo sea capaz de navegar de forma autónoma en entornos altamente dinámicos y congestionados, sino que lo haga garantizando la integridad biológica de la carga y la seguridad física de las personas. A diferencia de los sistemas de transporte industrial convencionales, nuestra propuesta pone el foco en la agilidad de maniobra mediante ruedas Mecanum, la seguridad activa basada en el ecosistema ROS 2 y la integración con la infraestructura del edificio inteligente.

A lo largo de este documento, se detallará el proceso de concepción del sistema siguiendo una estructura técnica rigurosa. En primer lugar, se realizará una revisión del estado del arte, analizando los referentes actuales en robótica de servicio sanitario. Posteriormente, se definirán los requerimientos funcionales y no funcionales que rigen el diseño, seguidos de una descripción exhaustiva de la arquitectura de hardware, la selección de materiales y los componentes electrónicos. Finalmente, se abordarán los aspectos de infraestructura, ciberseguridad y cumplimiento normativo, concluyendo con una propuesta de implementación práctica en un entorno simulado.
